% We reframe evaluation of agent context compression as a metric system over the compression target, rather than as a loose discussion of downstream performance. Let the compressed context state be defined by four components: \(Q=(D,R,P,O)\), where \(D\) denotes task-relevant density, \(R\) recoverability, \(P\) propagation of compression error, and \(O\) operational overhead. Under this view, a method is preferred when it maximizes useful preservation and recoverability while minimizing error spread and control cost.

% Equivalently, we can express the evaluation objective as \(J=\alpha D+\beta R-\gamma P-\delta O\), where the weights reflect the priorities of the target domain and agent setting. This formulation makes comparison more explicit: methods should not only be judged by final task success, but by how well they preserve information, support restoration, limit long-horizon failure, and remain efficient under budget constraints.

% Therefore, the objective \(J\) provides the comparison criterion for this section, linking benchmark results to compression quality under the domain-specific constraints \(\mathrm{Constraint}_i\).
\noindent\textbf{Outcome-Centric Evaluation.}
Benchmarks still judge context compression mainly by downstream task success. Success rate, with token count or compression ratio as an auxiliary proxy, remains the default protocol. This is convenient but inadequate: success entangles compression with planning and execution, while shorter prompts reveal little about whether task-critical evidence survives. Cross-paper comparison is further blurred by changes in backbone, task set, and context budget.

% \noindent\textbf{Current practice.}
% Current benchmarks treat context compression mainly as an ingredient of downstream task performance. Success rate, with token count or compression ratio as a secondary proxy, remains the default protocol. This is convenient but fundamentally insufficient: task success conflates compression with planning and execution, while shorter prompts say little about how much task-relevant information remains. Cross-paper comparison is further weakened by changes in agent backbone, task set, and context budget.

\noindent\textbf{The Missing Object.} The missing object is not the task outcome but the compressed state itself. For agents, four dimensions are decisive: \textit{information density}, \textit{recoverability}, \textit{error propagation}, and \textit{controller overhead}. Together, they ask whether compression preserves task-critical evidence, supports later restoration, avoids long-horizon failure amplification, and remains efficient enough to justify its control cost.

% \noindent\textbf{What is missing.} The missing object of evaluation is the compressed state itself. For agents, four dimensions are decisive: \textit{information density}, \textit{recoverability}, \textit{error propagation}, and \textit{controller overhead}. Together, they ask whether compression preserves task-critical evidence, supports later restoration, avoids long-horizon failure amplification, and remains efficient enough to justify its own control cost.

\noindent\textbf{From Outcomes to State Quality.} We therefore evaluate compressed states through \(Q=(D,R,P,O)\), where \(D\) denotes task-relevant density, \(R\) recoverability, \(P\) propagation risk, and \(O\) operational overhead, and summarize comparison through \(J=\alpha D+\beta R-\gamma P-\delta O\). The point is simple but consequential: compression should be judged as state quality, not only task outcome. This reframing ties benchmark outcomes to compression quality under domain-specific constraints \(\mathrm{Constraint}_i\).
% \noindent\textbf{A sharper framework.}  We therefore evaluate a compressed state through \(Q=(D,R,P,O)\), where \(D\) denotes task-relevant density, \(R\) recoverability, \(P\) propagation risk, and \(O\) operational overhead, and summarize comparison through \(J=\alpha D+\beta R-\gamma P-\delta O\). The core shift is conceptual: evaluation should treat compression as a state-quality problem, not merely an outcome-only scoring problem. This reframing links benchmark outcomes back to compression quality under domain-specific constraints \(\mathrm{Constraint}_i\).
Taken together, this reframing motivates a broader discussion of recoverable, faithful, and long-horizon context management, which we defer to Appendix~\ref{ap:extended_future_directions}.
% Taken together, these directions point toward a broader concluding discussion on recoverable, faithful, and long-horizon context management, which we defer to Appendix~\ref{ap:extended_future_directions}.

